\section{Conclusion}
\label{sec:conclusion}

We have computed the rectification of a graphene funnel and of an
electronic Tesla valve in the ballistic, viscous and plasmonic transport
limits, from dc to $3\THz$ and from the small-signal limit to the gate
overdrive. The Tesla valve does not rectify ballistically, rectifies in the
viscous liquid only beyond the drift velocity the Fermi liquid allows, and
is an order of magnitude below the funnel at terahertz frequencies. The
funnel's terahertz response, $0.07$ to $0.34\AW$ into a $75\ \Omega$
antenna, is the Dyakonov-Shur asymmetry of its terminations; its shape
contributes as a matched horn near a taper resonance and not as a diode.
The conversion efficiency follows a square law up to a saturation power
that scales as the square of the gate overdrive and peaks below one per
cent at every density and gate distance; contact resistance attacks the
asymmetry itself. The ballistic shape mechanism, driven at $1\THz$,
survives the transit time only when the potential drop is confined to the
neck and then reaches the same responsivity and no more.

Geometry, then, does not make a graphene terahertz rectenna. What it makes
is a square-law detector with a Johnson-limited noise-equivalent power of
tens of pW/Hz$^{1/2}$ at room temperature and of order $1\pWrtHz$ at low
density at $100$~K, if contacts of $100\ \Omega\um$ and the assumed
scattering times can be had. Two things would change the answer and are
outside the models used here: a rectification mechanism that is neither
ballistic, viscous nor plasmonic, of which the photothermoelectric effect
is the obvious example, and a material whose electron liquid sustains a
higher drift velocity relative to its plasma and Fermi velocities.
